% results_rq1

\section{RQ1: Gamification}
\label{sec:rq1}

A recurring claim in the self-tracking literature is that gamified
mechanics --- streaks, scores, badges --- induce anxiety and guilt in users who
fail to keep them up. Prayer trackers add a streak or a completion score to the
five daily prayers, so RQ1 asks whether this tension appears in devotional
tracking specifically. It does not, in the form the literature predicts.
Gamifying a prayer tracker neither helps nor harms sentiment, and guilt
language is \emph{rarer} inside tracker discourse than in the corpus at large.
What users resent is a record they cannot trust and a prompt that never
arrives. We present this as a reframing of the streak-anxiety hypothesis, not
merely a null result on it.

\subsection{Scope}

We define tracker discourse as the union of reviews carrying at least one of
three aspect tags --- \texttt{prayer\_tracker}, \texttt{tracker\_score},
\texttt{goal\_system} --- which yields 6,285 reviews across 25 of the 26
applications. The union is a definition rather than a measurement --- it is
declared as a three-element list in the released analysis code, and the tags it
names are keyword-defined in the released lexicon --- so what it needs is not
validation but a statement of what it contains. That composition is uneven.
At review level, \texttt{prayer\_tracker}, at 78\% precision, accounts for
5,907 of the 6,285; 378 reviews (6.0\% of the set) carry only
\texttt{tracker\_score} or \texttt{goal\_system}, both of which score 0\%
against their gold rows, the latter being the devotional-formula failure
described in \S\ref{sec:method-devotional-filtering}. These are review counts,
not the tagged-sentence counts in Table~\ref{tab:aspect-precision}, which are
not subtractable from them. The disparity is visible in the
definitions themselves: \texttt{prayer\_tracker} carries 13 keywords and 4
regular-expression patterns, against 2 keywords and 4 patterns for
\texttt{tracker\_score} and 7 keywords and 2 patterns for
\texttt{goal\_system}, so the two weaker tags rest on markedly less lexical
evidence. Those two aspects enter after the devotional-formula filter has already removed
the sentence pattern causing the \texttt{goal\_system} misfire. We report the two analyses that run on the full
set --- the tier comparison and the dip test --- on that basis, noting that
noise at this share attenuates a distributional signal rather than
manufacturing one. The guilt test and both coding passes draw their own
samples from within the set and do not depend on its outer boundary.

\subsection{No Sentiment Difference by Tracker Tier}

We compared sentiment across three tiers of tracking mechanic: no tracker,
a plain tracker, and a tracker combined with a numeric completion score. The
primary mixed-effects specification (M2) did not converge on this contrast ---
its two tier terms return estimates of 13.66 and 13.87 with standard errors of
\(2.0\times10^{6}\) and \(1.3\times10^{6}\) and \(p \approx 0.99999\) --- so we
report the tier comparison from the cluster-robust OLS sensitivity fit, which
did converge. Against no tracker, a plain tracker gives \(\beta = -0.087\)
(95\% CI \([-0.332, 0.158]\), \(p = 0.488\)) and a tracker with a numeric score
gives \(\beta = -0.134\) (\([-0.415, 0.147]\), \(p = 0.349\)), clustered on 25
apps (\(n = 6{,}285\)). Both intervals sit comfortably across zero. Adding a
streak or a numeric score to a prayer tracker does not measurably move how
users feel about the application.

\subsection{Tracker Sentiment Is Bimodal}

A null on the mean, however, was never going to be the right test: the
love/hate tension the literature predicts is a claim about the \emph{shape} of
the sentiment distribution, not its centre. Hartigan's dip test on tracker
valence confirms this shape directly: \(D = 0.118\) (\(n = 6{,}285\)), and the
reported significance underflows machine precision, so we report \(p <
10^{-16}\) rather than \(p = 0\). Figure~\ref{fig:rq1-bimodal} shows the
distribution: tracker sentiment splits toward the two poles rather than
clustering in the middle, with users loving the tracker or resenting it and
comparatively little ambivalence between. This is why the qualitative coding
below, rather than the tier comparison, carries the substantive answer to RQ1.

\begin{figure}[t]
  \centering
  \includegraphics[width=\linewidth]{fig06_rq1_tracker_sentiment}
  \caption{Distribution of tracker-aspect sentiment valence, split by application tracker tier. Hartigan's dip on the pooled tracker set gives \(D = 0.118\), \(p < 10^{-16}\).}
  \label{fig:rq1-bimodal}
\end{figure}

\subsection{Guilt Is Rarer Inside Tracker Discourse, Not Commoner}

If streak anxiety were driving the negative pole of that split, guilt language
should be over-represented in tracker discourse relative to the rest of the
corpus. Two independent lines of evidence say it is not, and we keep them
apart because they establish different things.

The first is lexical. We match two keyword lists over review text, one of
guilt- and anxiety-related terms and one of motivation-related terms. Within
tracker reviews the ratio of guilt to motivation matches is 0.26, against 0.86
elsewhere in the corpus (Fisher exact OR \(= 0.283\),
\(p = 3.5\times10^{-11}\), \(q = 6.7\times10^{-11}\), \(n = 149\))
(Figure~\ref{fig:rq1-guilt}). We state the ceiling on this plainly: unlike the
aspect tagger and the demand patterns, these two lists were never validated
against hand-labelled rows, and their vocabulary is not construct-pure ---
\emph{helpful} and \emph{improve} are ordinary praise terms, and \emph{fail}
and \emph{stress} are ordinary defect terms. The result is therefore a fact
about the relative frequency of two lexicons, not a validated measurement of
guilt, and we release the lists with the code so the reader can judge them.

The second line is the hand coding below, where two independently drawn samples
were read and labelled by a person rather than matched by pattern. It points
the same way, and it is the evidence we rest the construct claim on. Where guilt-adjacent language does appear inside
tracker reviews, it is frequently folded into an otherwise positive,
motivational account rather than expressing distress:

\begin{quote}
\itshape ``I love the app. My everyday friend who reminds me to pray, and makes
sure that I feel guilty if I don't [\normalfont sweat-smile emoji\itshape] I
recommend to everyone who wants to start praying, or keep the habit in a
consistent manner!''
\end{quote}
\hspace*{1em}--- Athan: Prayer Times \& Al Quran, 5 stars

\noindent We mark the emoji rather than drop it, since it is the clearest
signal in the review itself that the guilt described here is light-hearted.

\begin{figure}[t]
  \centering
  \includegraphics[width=\linewidth]{fig07_rq1_guilt_vs_motivation}
  \caption{Guilt- and motivation-coded affective language per application, as shares of that application's affect-coded reviews.}
  \label{fig:rq1-guilt}
\end{figure}

\subsection{Two Independent Coding Passes}
\label{sec:rq1-coding}

To characterise what tracker dissatisfaction is about, we coded two independent
samples of reviews expressing it --- this phrase rather than ``negative
reviews,'' since selection was by sentiment on the text rather than star
rating: 18 of the first sample's 50 carry 4--5 stars, so the set includes
otherwise-positive reviews built around one complaint. The first pass covered
50 reviews from a single application (\emph{Athan: Prayer Times \& Al Quran}),
selected before we identified that the selection rule's tie-break resolves
every three-valued sentiment tie in frame order, which here means app order;
the single-application sample is an artefact of that tie-break, not evidence
that one application dominates tracker complaints. The second pass drew 30
reviews stratified across the 22 other applications carrying codable tracker
dissatisfaction, no app contributing more than two, coded against the frame the
first pass produced.
Pass 1 yielded 33 open codes over 50 reviews (1 excluded, \(n = 49\)); pass 2
yielded 25 codes over 30 (2 excluded, \(n = 28\)) and required two new themes.
We report the passes side by side and never pool them, since their sampling
designs differ and a pooled percentage would silently re-weight the ecosystem
toward Athan (Table~\ref{tab:rq1-themes}).

\begin{table}[t]
  \centering
  \small
  \caption{Theme prevalence, coded independently in each pass. Never pooled.}
  \label{tab:rq1-themes}
  \begin{tabular}{lrr}
    \toprule
    Theme & Pass 1 (Athan, \(n{=}49\)) & Pass 2 (22 apps, \(n{=}28\)) \\
    \midrule
    C. Record can't be trusted        & 11 (22\%) & 6 (21\%) \\
    D. Normative mismatch             & 6 (12\%)  & 3 (11\%) \\
    E. Prompting failures             & 7 (14\%)  & \textbf{13 (46\%)} \\
    A. Redesign as regression         & 10 (20\%) & \textbf{0} \\
    B. Retroactive logging denied     & 5 (10\%)  & \textbf{0} \\
    F. Monetisation vs.\ religious purpose & 10 (20\%) & 1 (4\%) \\
    G. Record gated, not owned        & ---       & 3 (11\%)\(^\ast\) \\
    H. App overreaches into device    & ---       & 2 (7\%)\(^\ast\) \\
    \bottomrule
  \end{tabular}
  \\[2pt]
  {\footnotesize \(^\ast\)new theme, not present in the pass-1 frame}
\end{table}

Guilt or judgment is rare in both passes, and we give the two rates separately
rather than pooling them: 2 of the 49 coded complaints in the single-app pass
(4.1\%; \emph{tracker framing feels judgmental} and \emph{imposed goals replace
own tracking}) and 1 of the 28 in the cross-app pass (3.6\%; \emph{reminders
induce guilt}). A hypothesis this prominent in the self-tracking literature
surfacing at four percent or below in each of two independently drawn samples,
while the Fisher test above points the same way, is a robust null rather than
an underpowered one.

The two passes also demonstrate a general methodological hazard for app-store
review mining. Themes A, B, and F together account for 51\% of the single-app
pass's coded complaints but only 4\% of the cross-app pass's --- a concrete
case of single-app qualitative sampling manufacturing what looks like a
genre-level theme from what is, on a wider sample, a property of one
application's release.

\subsection{What Dissatisfaction Is Actually About}

Coding the cross-app sample against the shared frame surfaces a consistent
ordering, and it replicates across both independently drawn samples for two of
the three themes.

\textbf{Prompting failure is the dominant complaint ecosystem-wide (46\% of
the cross-app pass).} The tracker does not fail because logging is onerous;
it fails because the app never reliably prompts in the first place ---
location detection failing, adhan playback not firing, notifications lost,
calendars drifting out of sync:

\begin{quote}
\itshape ``the adhan is not working unless I open the app without logging out
of it till prayer time''
\end{quote}
\hspace*{1em}--- Pray Watch, 1 star

\textbf{Record trustworthiness replicates almost exactly across both samples}
(22\% in pass 1, 21\% in pass 2): users lose logged history, watch counts
regress, or find the tracker has silently disabled itself, often inside
reviews that otherwise praise the app:

\begin{quote}
\itshape ``I would give it a 5 star, but the prayers that I log on the
previous days or current day would sometimes disappear. Please fix this!''
\end{quote}
\hspace*{1em}--- Athan: Prayer Times \& Al Quran, 4 stars

\textbf{Normative mismatch is the smallest of the three replicating themes}
(\(\sim\)11\% in both passes) and is not the same construct as guilt. Where a
tracker review raises a genuinely normative objection, it is that the app
\emph{mis-models} religious practice, not that it makes the user feel
judged: menstrual exemption from prayer is unsupported, there is no support
for supererogatory (\emph{sunnah}) practice, and imposed goals displace the
user's own targets. We treat this distinction as the paper's contribution
here, not a restatement of the streak-anxiety story:

\begin{quote}
\itshape ``It's a really good app but I'm taking 1 star away cause I wish
there was some option for menstruating women cuz my overall score on prayer
tracker becomes really low when I can't pray for my period.''
\end{quote}
\hspace*{1em}--- Deen - Islamic App, 3 stars

\noindent Menstrual-exemption complaints of this kind appear independently in
\emph{both} coding passes, from different applications, and we return to them
as direct evidence for RQ4 (\S\ref{sec:rq4}).

\subsection{Themes A and B Are an Application-Specific Episode}

Two of the six original themes --- redesign as regression (20\% of pass 1) and
retroactive logging denied (10\%) --- vanish entirely in the cross-app pass.
Both concern a single redesign of the Athan tracker that broke an existing
logging workflow and removed retroactive logging. Because this pair accounted
for 30\% of the single-app sample's coded complaints, we report it as an
application-specific episode rather than folding it into any ecosystem-level
theme; conflating the two would overstate how much of tracker dissatisfaction
is about record integrity in general.

\subsection{Summary}

Adding a streak or a score does not move sentiment, and tracker sentiment is
bimodal, so a mean was never going to characterise it. Guilt language is
under-represented rather than over-represented inside tracker discourse, and
appears in 2 of 49 coded complaints in the single-app pass and 1 of 28 in the
cross-app pass. What tracker dissatisfaction is actually about, ecosystem-wide, is
the app failing to prompt reliably and failing to keep a trustworthy record;
where a normative objection surfaces, it concerns the app's model of worship
rather than pressure to perform.
